% FILE: chapters/ch08_s2_serving.tex
\section{Serving Models}\label{sec:serving_models}

In Section~\ref{sec:experiment_management}, we established rigorous practices for tracking machine learning experiments, versioning datasets with DVC, and registering validated model artifacts in an automated registry.

However, a trained model file sitting in cloud storage (such as a serialized \code{model.onnx}, \code{model.pt}, or \code{model.pkl}) delivers zero operational utility on its own. To extract real-world value, the model must be \textbf{served}: exposed as an always-on, network-accessible \textbf{Application Programming Interface (API)} that accepts telemetry from thousands of remote IoT devices, executes real-time inferences, and returns structured predictions.

Historically, serving a model was fraught with deployment friction. Machine learning engineers wrote scripts in research environments, and systems engineers struggled to deploy those scripts on enterprise servers due to mismatched Python versions, missing C libraries, and conflicting dependencies---the infamous \emph{``it worked on my laptop, but fails on the server!''} syndrome.

Modern Machine Learning Operations (MLOps) resolves this challenge through a standardized, three-stage serving architecture:
\begin{enumerate}
    \item \textbf{High-Performance API Framework (FastAPI)}: Wrapping the model inside an asynchronous, type-safe REST microservice that validates incoming sensor payloads.
    \item \textbf{Containerization (Docker)}: Packaging the model weights, application code, Python runtime, and system dependencies into an immutable, self-contained container image.
    \item \textbf{Cloud Deployment}: Running and scaling the containerized service on modern cloud infrastructure (such as serverless container platforms or managed container services).
\end{enumerate}

In this section, we study how to build production-ready ML serving APIs with FastAPI, master Docker containerization principles, plan cloud capacity mathematically using Little's Law, and trace the flow of edge telemetry into cloud-hosted models.

% ==============================================================================
% 1. BUILDING APIS WITH FASTAPI
% ==============================================================================
\subsection{Wrapping a Model in a FastAPI Service}\label{subsec:fastapi_service}

\paragraph{Why FastAPI Over Flask and Django?}
In the early days of Python web development, Flask and Django were the default choices for serving machine learning models. However, modern IoT telemetry pipelines demand high concurrency, minimal latency, and strict schema validation. \textbf{FastAPI} (created by Sebastián Ramírez) has become the de facto industry standard for ML model serving for several reasons:
\begin{itemize}
    \item \textbf{Asynchronous Non-Blocking I/O (ASGI)}: Built on Starlette, FastAPI natively supports Python's \code{async} and \code{await} syntax. A single server process can handle thousands of concurrent incoming HTTP connections while waiting for database queries or network sockets.
    \item \textbf{Automatic Data Validation via Pydantic}: Incoming sensor payloads are strictly validated against Python type hints. If an edge device transmits a corrupted string instead of a float, FastAPI intercepts the request instantly and returns a structured \code{422 Unprocessable Entity} error before the invalid data touches the model.
    \item \textbf{Automatic Interactive Documentation}: FastAPI automatically generates interactive OpenAPI (Swagger UI) documentation at \code{/docs}, allowing engineers to test endpoints interactively in a web browser without writing documentation by hand.
\end{itemize}

\begin{figure}[htbp]
\centering
\begin{tikzpicture}[
    scale=0.90,
    client/.style={rectangle, draw=keyideablue, fill=keyideablue!12, rounded corners=2mm, thick, minimum width=2.6cm, minimum height=1.3cm, align=center, font=\scriptsize},
    fastapi/.style={rectangle, draw=analogygreen, fill=analogygreen!12, rounded corners=3mm, thick, minimum width=6.8cm, minimum height=3.6cm, align=center},
    innerbox/.style={rectangle, draw=gray!70, fill=white, rounded corners=2mm, thick, minimum width=6.0cm, minimum height=0.9cm, align=center, font=\scriptsize},
    arrow/.style={-{Stealth[length=2.5mm]}, thick, draw=gray!80}
]
    % Client
    \node[client] (esp32) at (0, 0) {
        \textbf{Edge Device}\\
        (ESP32 / Gateway)\\
        \tiny Telemetry JSON Payload
    };

    % FastAPI Container
    \node[fastapi] (server) at (7.0, 0) {};
    \node[font=\small\bfseries, text=analogygreen!80!black] at (7.0, 1.4) {FastAPI Microservice (Uvicorn Server)};
    \node[innerbox] (pydantic) at (7.0, 0.6) {1. Pydantic Schema Validation (\code{TelemetryRequest})};
    \node[innerbox, draw=warningred, fill=warningred!10] (model) at (7.0, -0.4) {\textbf{2. In-Memory ML Model (\code{model.predict()})}};
    \node[innerbox] (response) at (7.0, -1.4) {3. JSON Response Serialization (\code{PredictionResponse})};

    % Network Flow
    \draw[arrow, draw=keyideablue] (esp32.15) -- node[above, font=\tiny\bfseries, text=keyideablue] {HTTP POST \code{/predict}} (pydantic.west);
    \draw[arrow, draw=analogygreen] (response.west) -- node[below, font=\tiny\bfseries, text=analogygreen] {HTTP 200 OK (Predictions)} (esp32.345);
\end{tikzpicture}
\caption{The request-validation-inference lifecycle inside a FastAPI model serving architecture.}\label{fig:fastapi_lifecycle}
\end{figure}

\paragraph{The Lifespan Pattern: Loading Models Once}
A catastrophic anti-pattern in model serving is loading model weights from disk \emph{inside} the request handling function:
\begin{lstlisting}[style=pythonstyle, numbers=none]
# BAD PRACTICE: NEVER DO THIS!
@app.post("/predict")
def predict(payload: TelemetryRequest):
    model = joblib.load("model.pkl")  # Loads 200 MB from disk on EVERY HTTP request!
    return {"prediction": model.predict(payload.features)}
\end{lstlisting}
Reading multi-megabyte weights from disk and rebuilding computational graphs on every incoming request increases response latency from $5\,\text{ms}$ to $500\,\text{ms}$ and quickly exhausts server memory under concurrent load!

The correct production approach is the \textbf{Lifespan Pattern}: model weights are loaded into RAM or GPU VRAM \textbf{exactly once} during application startup, and persist in memory across all incoming HTTP requests:
\begin{lstlisting}[style=pythonstyle, numbers=none]
from contextlib import asynccontextmanager
from fastapi import FastAPI
import joblib

ml_models = {}

@asynccontextmanager
async def lifespan(app: FastAPI):
    # Executed once before the server begins accepting HTTP requests
    print("[Startup] Loading model weights into system RAM...")
    ml_models["classifier"] = joblib.load("model.pkl")
    yield
    # Executed once when the server shuts down
    print("[Shutdown] Cleaning up model memory...")
    ml_models.clear()

app = FastAPI(lifespan=lifespan)
\end{lstlisting}

% ==============================================================================
% 2. DOCKER BASICS FOR MACHINE LEARNING
% ==============================================================================
\subsection{Docker Basics: Packaging the Machine Learning Runtime}\label{subsec:docker_basics}

\begin{analogy}[The Intermodal Shipping Container]
Before the 1950s, international shipping was slow, hazardous, and expensive. Cargo arrived at ports as loose sacks of grain, barrels of rum, and wooden crates of machinery. Longshoremen spent days manually loading and re-packing goods between horse carts, steamships, and trains. If a barrel broke, it ruined adjacent cargo.

In 1956, Malcolm McLean invented the \textbf{standardized intermodal shipping container}: a corrugated steel box of fixed dimensions ($20$ or $40$ feet). Once goods are packed inside, the container can be hoisted seamlessly onto a container ship, transferred to a freight train, or latched onto a semi-truck anywhere on Earth without ever opening the doors.

\textbf{Docker does for software what the shipping container did for global commerce.} Instead of shipping loose Python files and hoping the host server has the correct dependencies installed, Docker packages the application code, the exact Python interpreter, system shared libraries (\code{libc}, \code{libgomp}), environment variables, and model weights into an immutable, sealed software container that runs identically on any operating system.
\end{analogy}

\begin{figure}[htbp]
\centering
\begin{tikzpicture}[
    scale=0.88,
    box/.style={rectangle, draw=gray!70, fill=gray!10, rounded corners=2mm, thick, minimum width=5.5cm, minimum height=0.75cm, align=center, font=\scriptsize},
    vmbox/.style={rectangle, draw=keyideablue, fill=keyideablue!12, rounded corners=2mm, thick, minimum width=5.5cm, minimum height=0.75cm, align=center, font=\scriptsize},
    dockerbox/.style={rectangle, draw=analogygreen, fill=analogygreen!12, rounded corners=2mm, thick, minimum width=5.5cm, minimum height=0.75cm, align=center, font=\scriptsize}
]
    % Virtual Machine Architecture (Left)
    \node[font=\small\bfseries, text=keyideablue] at (2.75, 4.0) {Virtual Machines (VMs)};
    \node[vmbox] at (2.75, 3.2) {App A + Python Runtime};
    \node[vmbox, fill=keyideablue!20] at (2.75, 2.3) {Guest Operating System (Ubuntu, 4 GB)};
    \node[vmbox] at (2.75, 1.4) {Hypervisor (Type 1 or Type 2)};
    \node[box] at (2.75, 0.5) {Host Hardware (CPU, RAM, Disks)};

    % Docker Container Architecture (Right)
    \node[font=\small\bfseries, text=analogygreen] at (10.25, 4.0) {Docker Containers};
    \node[dockerbox] at (10.25, 3.2) {App A + Model (Isolated Namespaces)};
    \node[dockerbox, fill=analogygreen!20] at (10.25, 2.3) {Docker Engine / Container Runtime};
    \node[dockerbox] at (10.25, 1.4) {Host Linux Kernel (Shared cgroups)};
    \node[box] at (10.25, 0.5) {Host Hardware (CPU, RAM, Disks)};
\end{tikzpicture}
\caption{Architectural comparison: Virtual Machines require a complete guest OS per application, whereas Docker containers share the host Linux kernel, providing lightweight isolation.}\label{fig:vm_vs_docker}
\end{figure}

\paragraph{Containers vs.\ Virtual Machines (VMs)}
\begin{itemize}
    \item \textbf{Virtual Machines}: Emulate physical computer hardware using a hypervisor. Each VM runs its own full \textbf{guest operating system} (requiring gigabytes of storage, full system startup routines, and dedicated virtual memory allocations). Booting a VM takes minutes.
    \item \textbf{Containers}: Share the \textbf{host Linux kernel}. They achieve process isolation using native Linux kernel primitives:
    \begin{itemize}
        \item \textbf{Namespaces}: Isolate process trees (\code{pid}), network interfaces (\code{net}), and file mount points (\code{mnt}). A container believes it is running on a dedicated machine.
        \item \textbf{Control Groups (cgroups)}: Enforce hard quotas on hardware resources (e.g., restricting a container to exactly $2$ CPU cores and $1\,\text{GB}$ of RAM).
    \end{itemize}
    Because containers share the host kernel, they instantiate in \textbf{milliseconds} and consume negligible memory overhead compared to bare-metal execution.
\end{itemize}

\paragraph{Core Docker Terminology}
\begin{enumerate}
    \item \textbf{Dockerfile}: A plaintext recipe file containing sequential instructions used to assemble a container image.
    \item \textbf{Docker Image}: The read-only, immutable binary snapshot built from a Dockerfile. Images are composed of layered filesystems.
    \item \textbf{Docker Container}: A live, running, stateful instance of an image.
    \item \textbf{Docker Hub / Container Registry}: A centralized repository (e.g., AWS Elastic Container Registry, Google Artifact Registry) used to store and distribute versioned images across cloud servers.
\end{enumerate}

% ==============================================================================
% 3. CLOUD CAPACITY PLANNING MATHEMATICS
% ==============================================================================
\subsection{Cloud Capacity Planning and Concurrency}\label{subsec:capacity_planning}

Deploying a machine learning model to production requires sizing cloud server resources accurately. If too few server instances are provisioned, incoming requests queue up, latencies spike, and edge devices experience network timeouts. If too many instances are provisioned, the organization wastes thousands of dollars on idle compute.

\paragraph{Little's Law in Queuing Theory}
The relationship between request arrival rates, service latency, and concurrent requests in an ML serving pipeline is governed by \textbf{Little's Law} (John Little, 1961):
\begin{equation}
L = \lambda \cdot W\label{eq:littles_law}
\end{equation}
where:
\begin{itemize}
    \item $L$ is the average number of \textbf{concurrent requests} being processed inside the serving system.
    \item $\lambda$ is the \textbf{arrival rate} of incoming requests (requests per second, $\text{req/s}$).
    \item $W$ is the average \textbf{residence time} (the total latency required to process one request, in seconds).
\end{itemize}

\begin{example}[Hand Calculation: Sizing an IoT Cloud Prediction Cluster]
Let us calculate the exact server worker capacity required to handle an industrial fleet of connected edge sensors.

\paragraph{Given System Parameters:}
\begin{itemize}
    \item Fleet size: $6{,}000$ industrial motors, each equipped with an ESP32 edge gateway.
    \item Telemetry frequency: Each motor transmits an extracted statistical feature vector once every $20\,\text{seconds}$.
    \item Model inference latency: Profiling reveals that the cloud model takes an average of $t_{\text{inf}} = 25\,\text{ms}$ ($0.025\,\text{seconds}$) to process one feature vector.
    \item Target CPU utilization: Industry best practice dictates operating servers at no more than $70\%$ capacity ($U_{\text{target}} = 0.70$) to absorb sudden traffic bursts without queuing.
\end{itemize}

\paragraph{Step 1: Calculate Total Request Arrival Rate ($\lambda$)}
The total number of incoming requests arriving per second across all 6,000 devices is:
\begin{equation*}
\lambda = \frac{N_{\text{devices}}}{T_{\text{interval}}} = \frac{6{,}000 \text{ devices}}{20 \text{ seconds}} = \mathbf{300.0\,\text{requests/second}}
\end{equation*}

\paragraph{Step 2: Calculate Average Concurrent Requests ($L$)}
Using Little's Law (Equation~\eqref{eq:littles_law}):
\begin{equation*}
L = \lambda \times W = 300.0\,\text{req/s} \times 0.025\,\text{s} = \mathbf{7.5 \text{ concurrent requests}}
\end{equation*}
At any arbitrary millisecond in time, the cloud server cluster will actively be processing an average of $7.5$ requests simultaneously.

\paragraph{Step 3: Calculate Worker Throughput Capacity ($\mu$)}
A single synchronous worker thread executing requests with latency $W = 0.025\,\text{s}$ can process:
\begin{equation*}
\mu_{\text{single\_worker}} = \frac{1}{W} = \frac{1}{0.025\,\text{s}} = \mathbf{40.0\,\text{requests/second}}
\end{equation*}

\paragraph{Step 4: Determine Required Worker Threads}
To handle an arrival rate of $\lambda = 300.0\,\text{req/s}$ at $100\%$ saturation:
\begin{equation*}
N_{\text{workers\_nominal}} = \frac{\lambda}{\mu_{\text{single\_worker}}} = \frac{300.0}{40.0} = 7.5 \text{ workers}
\end{equation*}
Now, apply the $70\%$ utilization safety buffer ($U_{\text{target}} = 0.70$):
\begin{equation*}
N_{\text{workers\_safe}} = \frac{N_{\text{workers\_nominal}}}{U_{\text{target}}} = \frac{7.5}{0.70} \approx 10.71 \implies \mathbf{11 \text{ worker processes}}
\end{equation*}

\paragraph{Step 5: Translate to Container Sizing}
Standard production ASGI servers (such as Uvicorn managed by Gunicorn) recommend provisioning $2$ worker processes per CPU core:
\begin{equation*}
\text{Required CPU Cores} = \frac{N_{\text{workers\_safe}}}{2} = \frac{11}{2} = 5.5 \text{ cores} \implies \mathbf{6 \text{ vCPU cores}}
\end{equation*}
\textbf{Conclusion}: The engineering team should provision a cluster totaling \textbf{6 vCPUs} (for example, three 2-vCPU container instances running in a load-balanced serverless cluster). This guarantees that all $6{,}000$ industrial motors receive sub-$30\,\text{ms}$ diagnostic responses without request timeouts!
\end{example}

% ==============================================================================
% 4. CODE IMPLEMENTATION: FASTAPI SERVICE
% ==============================================================================
\subsection{Python Implementation: Production-Grade Model Server}\label{subsec:fastapi_code}

Listing~\ref{lst:fastapi_server} demonstrates a complete, production-grade model serving script using FastAPI, Pydantic, and Scikit-Learn.

\begin{lstlisting}[style=pythonstyle, caption={Complete FastAPI model serving application with Pydantic validation.}, label={lst:fastapi_server}]
import time
from contextlib import asynccontextmanager
from typing import List
from fastapi import FastAPI, HTTPException, status
from pydantic import BaseModel, Field
import numpy as np
from sklearn.ensemble import RandomForestClassifier

# 1. Define Strict Pydantic Input/Output Schemas
class SensorTelemetryRequest(BaseModel):
    device_id: str = Field(..., example="ESP32_PUMP_04", description="Unique edge identifier")
    timestamp: int = Field(..., example=1700000000, description="Unix epoch timestamp")
    features: List[float] = Field(
        ..., 
        min_items=4, 
        max_items=4, 
        example=[0.42, 1.85, 3.12, 4.40],
        description="[RMS, Peak-to-Peak, Kurtosis, CrestFactor]"
    )

class DiagnosticResponse(BaseModel):
    device_id: str
    status: str
    anomaly_probability: float
    execution_time_ms: float
    action_required: bool

# 2. Global In-Memory Model State
model_store = {}

# 3. Lifespan Event: Pre-load Model Weights Once
@asynccontextmanager
async def lifespan(app: FastAPI):
    print("[Startup] Initializing machine learning model...")
    # Simulate a pre-trained Random Forest model
    clf = RandomForestClassifier(n_estimators=20, random_state=42)
    X_mock = np.random.randn(200, 4)
    y_mock = (X_mock[:, 1] > 1.0).astype(int)
    clf.fit(X_mock, y_mock)
    
    model_store["classifier"] = clf
    print("[Startup] Model weights loaded into memory successfully.")
    yield
    print("[Shutdown] Releasing model resources.")
    model_store.clear()

# 4. Instantiate FastAPI Application
app = FastAPI(
    title="Industrial Condition Monitoring API",
    description="Real-time predictive maintenance serving endpoint for edge IoT nodes",
    version="1.0.0",
    lifespan=lifespan
)

# 5. Health Check Endpoints (For Cloud Load Balancers)
@app.get("/healthz", status_code=status.HTTP_200_OK)
def health_check():
    if "classifier" not in model_store:
        raise HTTPException(status_code=503, detail="Model not initialized")
    return {"status": "HEALTHY", "model_loaded": True}

# 6. Core Prediction Route
@app.post("/v1/predict", response_model=DiagnosticResponse)
async def predict_machine_health(payload: SensorTelemetryRequest):
    t_start = time.perf_counter()
    
    classifier = model_store.get("classifier")
    if not classifier:
        raise HTTPException(status_code=500, detail="Inference engine unavailable")
    
    # Reshape input into 2D NumPy array: (1, 4)
    input_data = np.array(payload.features).reshape(1, -1)
    
    # Execute Model Inference
    try:
        prediction = int(classifier.predict(input_data)[0])
        probabilities = classifier.predict_proba(input_data)[0]
        anomaly_prob = float(probabilities[1])
    except Exception as e:
        raise HTTPException(status_code=400, detail=f"Inference error: {str(e)}")
        
    t_latency = (time.perf_counter() - t_start) * 1000.0  # Convert to ms
    
    health_status = "CRITICAL_FAULT" if prediction == 1 else "NORMAL"
    return DiagnosticResponse(
        device_id=payload.device_id,
        status=health_status,
        anomaly_probability=round(anomaly_prob, 4),
        execution_time_ms=round(t_latency, 2),
        action_required=(prediction == 1)
    )
\end{lstlisting}

\paragraph{Line-by-Line Code Explanation}
\begin{itemize}
    \item \textbf{Lines 9--19}: Defines the \code{SensorTelemetryRequest} schema using Pydantic. Notice \code{min_items=4, max_items=4}: if an edge node sends 3 or 5 numbers, Pydantic immediately rejects the request with an HTTP 422 error, shielding the model from array shape crashes.
    \item \textbf{Lines 29--41}: The \code{lifespan} context manager loads model weights during application boot, guaranteeing that weights are stored in RAM permanently.
    \item \textbf{Lines 53--57}: The \code{/healthz} endpoint is used by cloud load balancers (e.g., AWS Application Load Balancer or Kubernetes liveness probes) to verify that the container is alive and ready to accept traffic.
    \item \textbf{Lines 60--88}: The \code{/v1/predict} route receives the validated payload, executes inference via \code{classifier.predict_proba()}, logs execution latency in milliseconds, and returns a structured response matching \code{DiagnosticResponse}.
\end{itemize}

% ==============================================================================
% 5. DOCKERFILE IMPLEMENTATION
% ==============================================================================
\subsection{Dockerizing the Model Server}\label{subsec:dockerfile}

Listing~\ref{lst:dockerfile} presents a production-grade, multi-stage \textbf{Dockerfile}. It adheres to cloud-native security best practices: pinning base images, minimizing layer counts, and running as a non-privileged user.

\begin{lstlisting}[style=pythonstyle, caption={Production multi-stage Dockerfile for FastAPI model serving.}, label={lst:dockerfile}]
# ==============================================================================
# Stage 1: Build & Dependencies
# ==============================================================================
FROM python:3.11-slim as builder

WORKDIR /install

# Install build-time compiler dependencies
RUN apt-get update && apt-get install -y --no-install-recommends \
    build-essential \
    && rm -rf /var/lib/apt/lists/*

# Copy and install Python requirements to /install directory
COPY requirements.txt .
RUN pip install --no-cache-dir --prefix=/install -r requirements.txt

# ==============================================================================
# Stage 2: Final Minimal Runtime Image
# ==============================================================================
FROM python:3.11-slim as runtime

WORKDIR /app

# Copy only the compiled Python packages from the builder stage
COPY --from=builder /install /usr/local

# Copy application source code and model weights
COPY main.py .
COPY model.pkl .

# Security Best Practice: Create and switch to a non-root system user
RUN useradd -m -u 1000 appuser && chown -R appuser:appuser /app
USER appuser

# Expose HTTP port 8000
EXPOSE 8000

# Start Uvicorn production server (4 worker processes on 0.0.0.0:8000)
CMD ["uvicorn", "main:app", "--host", "0.0.0.0", "--port", "8000", "--workers", "4"]
\end{lstlisting}

\paragraph{Building and Running the Container}
To build and execute the containerized service locally, run:
\begin{lstlisting}[style=pythonstyle, numbers=none]
# 1. Build the Docker image tagged as "tinyml-server:v1"
docker build -t tinyml-server:v1 .

# 2. Run container: map host port 8000 to container port 8000
docker run -d -p 8000:8000 --name edge_ml_api tinyml-server:v1

# 3. Test endpoint using curl from an edge device terminal:
curl -X POST "http://localhost:8000/v1/predict" \
     -H "Content-Type: application/json" \
     -d '{"device_id":"ESP32_01","timestamp":1700000000,"features":[0.42,1.85,3.12,4.40]}'
\end{lstlisting}

\begin{lstlisting}[style=outputstyle, caption={HTTP JSON response returned from the containerized FastAPI server.}]
{
  "device_id": "ESP32_01",
  "status": "NORMAL",
  "anomaly_probability": 0.0450,
  "execution_time_ms": 1.84,
  "action_required": false
}
\end{lstlisting}

% ==============================================================================
% 6. DEPLOYING ON CLOUD SERVICES
% ==============================================================================
\subsection{Deploying to Cloud Services}\label{subsec:cloud_deployment}

Once an application is packaged into a Docker container, it can be deployed to any cloud provider with zero code modifications:

\paragraph{1. Serverless Containers (Google Cloud Run / AWS App Runner)}
For IoT telemetry systems, \textbf{serverless container platforms} offer the ideal balance of cost and operational simplicity:
\begin{itemize}
    \item \textbf{Scale to Zero}: When no IoT sensors are transmitting data (e.g., during off-hours in an industrial plant), the cloud provider shuts down all container instances, dropping hosting costs to \textbf{\$0.00}.
    \item \textbf{Automatic Scaling}: As sensor fleets wake up, the platform detects incoming HTTP traffic and automatically spins up dozens of container replicas within milliseconds.
    \item \textbf{Deployment Command (Google Cloud Run)}:
    \begin{lstlisting}[style=pythonstyle, numbers=none]
    gcloud run deploy industrial-ml-service \
        --image gcr.io/my-project/tinyml-server:v1 \
        --platform managed \
        --region us-central1 \
        --memory 1Gi \
        --min-instances 0 \
        --max-instances 20
    \end{lstlisting}
\end{itemize}

\paragraph{2. Managed Kubernetes (EKS / GKE)}
For massive enterprise fleets with hundreds of thousands of concurrent sensors, organizations deploy container images onto \textbf{Kubernetes}:
\begin{itemize}
    \item \textbf{Horizontal Pod Autoscaling (HPA)}: Dynamically scales container pods based on CPU utilization or incoming HTTP request queues.
    \item \textbf{Rolling Deployments}: Updates model versions with zero downtime by spinning up containers with model v2.0 before gracefully terminating model v1.0 containers.
\end{itemize}

% ==============================================================================
% 7. IOT AND TINYML CONNECTION
% ==============================================================================
\begin{iotbox}[The Cascaded Edge-Cloud Serving Pipeline]
How does a cloud FastAPI microservice integrate with real-world edge IoT devices?

\paragraph{The Cascaded Architecture in Practice:}
\begin{enumerate}
    \item \textbf{Tier 1: On-Device TinyML Screening}: An industrial water pump runs an ultra-low-power Int8 model on an ESP32. The microcontroller evaluates vibration continuously at $1\,\text{kHz}$ using only $40\,\text{mA}$. For $99.9\%$ of operational cycles, the vibration is normal, and the ESP32 keeps its Wi-Fi radio off.
    \item \textbf{Tier 2: Cloud Escalation}: When local vibration exceeds an initial anomaly threshold, the ESP32 activates its Wi-Fi radio, formats a JSON payload containing the last 2 seconds of statistical features, and sends an \code{HTTP POST} to the cloud FastAPI endpoint hosted on Google Cloud Run.
    \item \textbf{Tier 3: Deep Cloud Diagnostics}: The cloud FastAPI service runs a massive ensemble model (or queries an LLM diagnostic engine), correlating the motor's vibration with ambient weather logs and maintenance databases. The cloud returns an advisory message: \code{"Bearing race wear: Schedule grease lubrication within 48 hours"}.
    \item \textbf{Result}: The edge device saves $99\%$ of its battery energy by remaining offline most of the time, while the cloud server only consumes compute resources when physical anomalies genuinely require deep analysis!
\end{enumerate}
\end{iotbox}

% ==============================================================================
% 8. COMMON MISTAKES AND MISCONCEPTIONS
% ==============================================================================
\begin{warning}[Common Serving and Containerization Mistakes]
\begin{itemize}
    \item \textbf{Running Uvicorn with \code{--reload} in Production}:
    The \code{--reload} flag tells Uvicorn to monitor filesystem changes and reload code automatically. This adds massive file-polling overhead, consumes CPU cycles, and can cause random worker restarts during live inference. Use \code{--reload} strictly during local development; never include it in a Dockerfile \code{CMD}.
    \item \textbf{Running Containers as Root}:
    By default, processes inside a Docker container execute with root privileges (\code{uid=0}). If an attacker exploits a remote code execution vulnerability in your FastAPI application or an underlying dependency, they gain root access to the container and potentially the host system! Always define an unprivileged user (\code{USER appuser}) in production Dockerfiles.
    \item \textbf{Neglecting Cold-Start Latency in Serverless Platforms}:
    When a serverless container scales from zero, the cloud provider must pull the image, allocate memory, boot the container, and initialize the model weights. If your PyTorch Docker image is $4\,\text{GB}$, this cold start can take $15$ to $30\,\text{seconds}$. An edge device with a 5-second HTTP timeout will abort and drop the alert! Keep Docker images small ($< 500\,\text{MB}$) by using multi-stage builds and \code{python:slim} base images.
    \item \textbf{Omitting Input Array Shape Validation}:
    Passing a 1D array of shape \code{(4,)} into a scikit-learn model that expects a 2D array of shape \code{(1, 4)} triggers an unhandled Python exception, returning a generic \code{500 Internal Server Error} to the IoT client. Always reshape arrays explicitly and wrap inference in \code{try-except} blocks.
\end{itemize}
\end{warning}

% ==============================================================================
% 9. KEY TAKEAWAYS
% ==============================================================================
\begin{takeaways}
\begin{itemize}
    \item \textbf{Production API Serving}: FastAPI provides an asynchronous, type-safe serving layer for machine learning models, validating incoming sensor data via Pydantic schemas before executing inference.
    \item \textbf{The Lifespan Pattern}: Model weights must be loaded into memory once during application startup, never dynamically re-read from disk inside individual HTTP request handlers.
    \item \textbf{Containerization with Docker}: Docker bundles code, dependencies, and model weights into an immutable container image, eliminating environmental drift between developer laptops and cloud servers.
    \item \textbf{Capacity Planning}: Little's Law ($L = \lambda W$) provides the mathematical foundation for sizing cloud compute clusters, determining the exact concurrency and worker threads needed to handle edge sensor arrival rates.
    \item \textbf{Cascaded Edge-Cloud Synergy}: On-device TinyML acts as a first-line filter, escalating only ambiguous or anomalous signals to cloud-hosted FastAPI microservices.
\end{itemize}
\end{takeaways}

% ==============================================================================
% 10. CHECK YOURSELF
% ==============================================================================
\begin{checkyourself}
\begin{enumerate}
    \item \textbf{Conceptual}: Why does loading machine learning model weights inside a FastAPI endpoint function (e.g., inside \code{def predict()}) cause severe production failures under concurrent load?
    \item \textbf{Calculation}: A cloud ML microservice receives an average arrival rate of $\lambda = 250\,\text{requests/second}$ from smart electric meters. Benchmarking reveals the model takes $40\,\text{ms}$ ($0.040\,\text{seconds}$) to process one request.
    \begin{enumerate}
        \item Using Little's Law ($L = \lambda W$), calculate the average number of concurrent requests $L$ actively being processed in the system.
        \item If a single worker process handles $25\,\text{requests/second}$, how many worker processes are required to operate at a safe $70\%$ capacity utilization?
    \end{enumerate}
    \item \textbf{Docker Security}: In a production Dockerfile, why is it considered dangerous to omit the \code{USER} directive, and what security principle does defining a non-root user satisfy?
    \item \textbf{Conceptual}: What is the primary difference between a Virtual Machine (VM) and a Docker Container in terms of kernel utilization and resource overhead?
\end{enumerate}
\end{checkyourself}

% ==============================================================================
% 11. ANSWERS
% ==============================================================================
\subsection*{Answers}

\begin{enumerate}
    \item \textbf{Consequences of Per-Request Model Loading}:
    Loading model weights from disk inside the request handler forces the operating system to perform disk I/O, parse file structures, allocate memory buffers, and rebuild the computational graph on \emph{every single incoming HTTP request}. Under concurrent traffic (e.g., 50 requests per second), this:
    \begin{itemize}
        \item Increases response latency by $10\times$ to $100\times$ (from milliseconds to seconds).
        \item Sinks CPU resources into repeated file decoding.
        \item Quickly exhausts system RAM, triggering Out-Of-Memory (OOM) crashes as dozens of redundant copies of the model are instantiated simultaneously.
    \end{itemize}
    The \textbf{Lifespan pattern} loads the model once at startup, keeping a single copy in memory to serve all subsequent requests.

    \item \textbf{Step-by-Step Capacity Planning Calculation}:
    Given:
    \begin{itemize}
        \item Arrival rate: $\lambda = 250\,\text{req/s}$
        \item Service latency: $W = 0.040\,\text{s}$
        \item Single worker capacity: $\mu = 25\,\text{req/s}$
        \item Target utilization: $U = 0.70$
    \end{itemize}
    \begin{enumerate}
        \item \textbf{Average Concurrent Requests ($L$)}:
        \begin{equation*}
        L = \lambda \times W = 250\,\text{req/s} \times 0.040\,\text{s} = \mathbf{10.0 \text{ concurrent requests}}
        \end{equation*}
        At any given moment, the cluster is actively computing exactly 10 requests.
        \item \textbf{Required Workers at 70\% Utilization}:
        First, compute nominal workers needed at $100\%$ capacity:
        \begin{equation*}
        N_{\text{nominal}} = \frac{\lambda}{\mu} = \frac{250}{25} = 10 \text{ workers}
        \end{equation*}
        Apply the $70\%$ safety utilization buffer:
        \begin{equation*}
        N_{\text{safe}} = \frac{N_{\text{nominal}}}{U} = \frac{10}{0.70} \approx 14.28 \implies \mathbf{15 \text{ worker processes}}
        \end{equation*}
        The team must provision at least 15 worker processes across the cluster.
    \end{enumerate}

    \item \textbf{Docker Non-Root Security Best Practice}:
    By default, processes inside a Docker container run as \textbf{root} (\code{uid=0}). If an attacker discovers an exploit in the web framework, a dependency library, or the Python runtime that allows arbitrary code execution, they obtain root privileges inside the container. In many container breakout vulnerabilities, container root privileges can be leveraged to compromise the underlying host server. Creating and specifying a dedicated, unprivileged user (e.g., \code{USER appuser}) satisfies the \textbf{Principle of Least Privilege}: even if the web service is compromised, the attacker cannot modify system files, install rootkits, or control host resources.

    \item \textbf{Virtual Machines vs.\ Docker Containers}:
    \begin{itemize}
        \item \textbf{Virtual Machines}: Run on top of a hypervisor that emulates physical hardware. Each VM requires its own complete \textbf{guest operating system} with a dedicated kernel, device drivers, and system background daemons. This results in heavy disk footprints ($5 - 20\,\text{GB}$) and high RAM overhead.
        \item \textbf{Docker Containers}: Share the \textbf{host Linux kernel}. They do not boot a separate operating system; instead, they isolate application processes using native kernel namespaces and cgroups. This allows containers to boot in milliseconds and operate with near-zero virtualization overhead.
    \end{itemize}
\end{enumerate}
